%%% FIELD fixed-gap-frontier-proof
We prove the statement with the displayed lower bound on the horizon.  First fix
\(P\in\mathcal P_{\Delta_0,b_0}^{\mathrm{fg}}\).  By the H\"older-means
assumption, \(\tau_P=\mu_1-\mu_0\) is continuous.  The fixed-gap assumption
holds \(P_X\)-almost everywhere.  If \(|\tau_P(x)|<\Delta_0\) at some
\(x\in[0,1]^d\), continuity gives a relatively open neighborhood \(G\) of
\(x\) on which \(|\tau_P|<\Delta_0\).  Such a neighborhood has positive
Lebesgue measure, and the context-density assumption gives
\[
 P_X(G)=\int_G p_X(u)\,du\geq c_X\operatorname{Leb}(G)>0,
\]
contradicting the almost-everywhere fixed-gap condition.  Hence
\(|\tau_P(x)|\geq\Delta_0\) everywhere.  If \(\tau_P\) took both signs, its
restriction to the line segment joining two points of opposite sign would be
continuous and would vanish at an intermediate point.  This is impossible.
Thus exactly one arm is superior for every context.

We now specify constants that make the pilot construction valid.  Put
\[
 v_0=4(M^2+\sigma^2),\qquad
 a_0=\frac{2v_0}{\Delta_0^2},\qquad
 K=\max\left\{4,\frac{2}{\delta},16M\right\},
\]
and
\[
 B=a_0\left(1+\frac{\log K}{\log 2}\right)+\frac{1}{\log 2},
 \qquad
 C_1=\max\left\{8MB,\frac{1}{\log 2}\right\},
 \qquad c_1=\frac12.
\]
Let \(D_0=\max\{1,B,C_1\}\) and
\[
 T_0=\left\lceil4D_0\log(4D_0)\right\rceil.
\]
For \(x_0=4D_0\log(4D_0)\), the inequalities
\(\log(4D_0)\leq4D_0\) and
\(x_0\leq(4D_0)^2\) give
\(D_0\log x_0\leq x_0/2\).  Since \(x/\log x\) is increasing for
\(x>e\), it follows that
\[
 D_0\log T\leq\frac{T}{2}
 \quad\text{for every integer }T\geq T_0.
\]
This is the promised explicit large-horizon restriction.

Fix henceforth \(T\geq T_0\) and
\(C_1\log T\leq\rho\leq c_1T\).  Define
\[
 k=\left\lceil a_0\log(KT)\right\rceil.
\]
Because \(T\geq2\),
\[
 k\leq a_0\left(1+\frac{\log K}{\log2}\right)\log T+1
 \leq B\log T\leq\frac{T}{2}.
\]
Moreover,
\[
 2Mk\leq2MB\log T\leq\frac{C_1\log T}{4}\leq\frac\rho4,
 \qquad
 \rho\geq C_1\log T\geq1.
\]

During rounds \(1,\ldots,k\), independently assign either arm with
probability \(1/2\), and set
\[
 Z_t=2(2A_t-1)Y_t,\qquad
 \widehat s=\operatorname{sign}\left(k^{-1}\sum_{t=1}^k Z_t\right),
\]
using a fixed deterministic rule at zero.  Let
\(\bar\tau_P=\mathbb E_{P_X}\tau_P(X)\).  Conditional on the pre-round
history, balanced assignment, independent contexts, and the Gaussian reward
model give
\[
 \mathbb E_{P,\mathsf D}[Z_t\mid H_{t-1}]=\bar\tau_P.
\]
The preceding global-sign argument gives
\(|\bar\tau_P|\geq\Delta_0\), with the sign of \(\bar\tau_P\) equal to the
superior-arm sign.

For completeness, we derive the conditional exponential bound used for the
pilot.  Write
\[
 S_t=2(2A_t-1)\mu_{A_t}(X_t),\qquad
 G_t=2(2A_t-1)\sigma\varepsilon_t,\qquad Z_t=S_t+G_t.
\]
The bounded-means assumption gives \(-2M\leq S_t\leq2M\).  If a centered
random variable \(V\) lies in an interval \([a,b]\), convexity gives
\[
 e^{\lambda V}
 \leq\frac{b-V}{b-a}e^{\lambda a}
     +\frac{V-a}{b-a}e^{\lambda b}.
\]
After expectation, put \(q=-a/(b-a)\).  The logarithm of the resulting upper
bound is
\[
 f(s)=\log\left\{(1-q)e^{-qs}+q e^{(1-q)s}\right\},
 \qquad s=\lambda(b-a).
\]
Here \(f(0)=f'(0)=0\), and direct differentiation shows that
\(f''(s)=q_s(1-q_s)\leq1/4\), where \(q_s\in[0,1]\) is the exponentially
tilted Bernoulli probability.  Integrating this second-derivative bound from
zero yields \(f(s)\leq s^2/8\) for either sign of \(s\).  Applying this
calculation conditionally to
\(S_t-\mathbb E[S_t\mid H_{t-1}]\), whose supporting interval has length
\(4M\), gives
\[
 \mathbb E\left[
 e^{\lambda(S_t-\bar\tau_P)}\mid H_{t-1}
 \right]\leq e^{2M^2\lambda^2}.
\]
By the Gaussian-innovation assumption,
\[
 \mathbb E[e^{\lambda G_t}\mid H_{t-1},X_t,A_t]
 =e^{2\sigma^2\lambda^2}.
\]
Iterated conditioning therefore gives, for every \(\lambda\in\mathbb R\),
\[
 \mathbb E\left[
 e^{\lambda(Z_t-\bar\tau_P)}\mid H_{t-1}
 \right]
 \leq e^{v_0\lambda^2/2}.
\]
Applying this inequality successively over the \(k\) pilot rounds yields
\[
 \mathbb E\exp\left\{
 \lambda\sum_{t=1}^k(Z_t-\bar\tau_P)
 \right\}
 \leq e^{kv_0\lambda^2/2}.
\]
Chernoff's argument, optimized at
\(|\bar\tau_P|/v_0\) and applied to the appropriate tail, now gives
\[
 \mathbb P_{P,\mathsf D}^{T}\{\widehat s\text{ is wrong}\}
 \leq
 \exp\left\{-\frac{k\bar\tau_P^2}{2v_0}\right\}
 \leq
 \exp\left\{-\frac{k\Delta_0^2}{2v_0}\right\}
 \leq\frac{1}{KT}.
\]

Let \(N=T-k\), so \(N\geq T/2\), and set
\[
 p=\min\left\{\frac14,\frac{\rho}{16MN}\right\}.
\]
In every remaining round, use fresh randomization to assign the arm declared
inferior with probability \(p\) and the arm declared superior with probability
\(1-p\).  This is a single non-anticipating design fixed independently of
\(P\).  The pilot loss is at most \(2Mk\leq\rho/4\).  On correct pilot
classification, the expected tracking loss is bounded by
\[
 2MpN\leq\frac\rho8.
\]
On incorrect classification, its conditional loss is at most \(2MN\), so its
unconditional contribution is at most
\[
 \frac{2MN}{KT}\leq\frac{2M}{K}\leq\frac18\leq\frac\rho8.
\]
Consequently, by the definition of participant-welfare loss,
\[
 \sup_{P\in\mathcal P_{\Delta_0,b_0}^{\mathrm{fg}}}
 \mathcal R_T(P,\mathsf D)
 \leq\frac\rho4+\frac\rho8+\frac\rho8
 =\frac\rho2\leq\rho.
\]
Thus \(\mathsf D\in\mathfrak D_T(\rho;
\mathcal P_{\Delta_0,b_0}^{\mathrm{fg}})\).  In addition, a split into the
two cases defining \(p\) gives
\[
 Np\geq c_M\rho,
 \qquad
 c_M=\min\left\{\frac{1}{16M},\frac14\right\}>0.
\]
Indeed, equality with \(\rho/(16M)\) holds when the second branch is active;
otherwise \(Np=N/4\geq T/8\geq\rho/4\), using
\(\rho\leq T/2\).

We next define one connected observed-data interval.  Use only the tracking
observations.  If \(\widehat s=+1\), let \(g(x)=1-\pi_0(x)\) and
\[
 U_t=g(X_t)\left\{
 \frac{A_tY_t}{1-p}-\frac{(1-A_t)Y_t}{p}
 \right\}.
\]
If \(\widehat s=-1\), let \(g(x)=\pi_0(x)\) and
\[
 U_t=g(X_t)\left\{
 \frac{(1-A_t)Y_t}{1-p}-\frac{A_tY_t}{p}
 \right\}.
\]
Conditional on the pilot, inverse-propensity cancellation gives
\[
 \mathbb E[U_t\mid\text{pilot}]
 =\begin{cases}
 \mathbb E_{P_X}[(1-\pi_0(X))\tau_P(X)],&\widehat s=+1,\\
 \mathbb E_{P_X}[\pi_0(X)(-\tau_P(X))],&\widehat s=-1.
 \end{cases}
\]
On correct pilot classification, the applicable quantity is exactly
\(\theta(P)=\mathbb E_{P_X}[\tau_P(X)_+-\pi_0(X)\tau_P(X)]\), by the
targeting-gain definition and the global-sign property.  Since
\(0\leq g\leq1\) and the Gaussian reward and bounded-means assumptions imply
\(\mathbb E[Y_t^2\mid X_t,A_t]\leq M^2+\sigma^2\), either sign satisfies
\[
 \mathbb E[U_t^2\mid\text{pilot}]
 \leq(M^2+\sigma^2)\left\{\frac1p+\frac1{1-p}\right\}
 \leq\frac{4(M^2+\sigma^2)}{3p},
\]
where the last inequality uses \(p\leq1/4\).  The tracking observations are
conditionally independent given the pilot, so the conditional variance of
their average is at most \(4(M^2+\sigma^2)/(3Np)\).

Define
\[
 \widehat\theta_T=\frac1N\sum_{t=k+1}^T U_t,
 \qquad
 r_T=\sqrt{\frac{8(M^2+\sigma^2)}{3\delta Np}},
 \qquad
 I_T=[\widehat\theta_T-r_T,\widehat\theta_T+r_T].
\]
The deterministic tie rule and measurability of \(\pi_0\) make both endpoints
Borel functions of \(O_{1:T}\); thus \(I_T\) is one \(P\)-independent
connected interval map.  Conditional Chebyshev gives
\[
 \mathbb P\left\{
 |\widehat\theta_T-\mathbb E[\widehat\theta_T\mid\text{pilot}]|>r_T
 \mathrel{\big|}\text{pilot}
 \right\}\leq\frac\delta2.
\]
Since \(1/(KT)\leq\delta/2\), the union bound with the pilot error yields
\[
 \inf_{P\in\mathcal P_{\Delta_0,b_0}^{\mathrm{fg}}}
 \mathbb P_{P,\mathsf D}^{T}\{\theta(P)\in I_T\}
 \geq1-\delta.
\]
Therefore
\(I_T\in\mathfrak C_{T,\delta}(\mathsf D;
\mathcal P_{\Delta_0,b_0}^{\mathrm{fg}})\).  Its length is deterministic and
obeys
\[
 |I_T|=2r_T
 \leq
 2\sqrt{\frac{8(M^2+\sigma^2)}{3\delta c_M}}\,\rho^{-1/2}.
\]
The definitions of the width envelope and budgeted designs prove the upper
bound.

It remains to prove the procedure-free converse.  The interiority condition
permits a constant \(\gamma\in(\Delta_0,2M)\) with the following margin
property.  If \(\alpha=0\), choose any such \(\gamma\); the
gap--margin-compatibility assumption gives \(C_0\geq1\).  If \(\alpha>0\),
choose
\[
 \max\{\Delta_0,\min(\delta_0,C_0^{-1/\alpha})\}<\gamma<2M.
\]
Then either \(\gamma>\delta_0\), or
\(C_0\gamma^\alpha>1\).  Thus a constant contrast of magnitude at least
\(\gamma\) satisfies the local margin assumption: for \(u\) below the
constant gap the relevant event is empty, while for \(u\) at or above the gap
its probability is one and the asserted inequality follows from
\(C_0u^\alpha\geq1\).  In the case \(\alpha=0\), the event at \(u=0\) is
empty and the same conclusion for positive \(u\) follows from \(C_0\geq1\).

Because the cube has unit volume,
\[
 \int_{\mathcal X}(1-\pi_0(x))\,dx+
 \int_{\mathcal X}\pi_0(x)\,dx=1.
\]
Choose arm one to be superior if the first integral is at least \(1/2\), and
choose arm zero to be superior otherwise.  Let \(a_{\mathrm{inf}}\) denote
the resulting inferior arm and let \(b\) denote the corresponding benchmark
mass.  Then \(b\geq1/2\geq b_0\).

Let \(P\) have uniform contexts, conditionally independent Gaussian outcomes
of variance \(\sigma^2\), superior-arm mean \(\gamma/2\), and inferior-arm
mean \(-\gamma/2\).  Choose a constant
\[
 0<c_*\leq\min\left\{\frac{2M-\gamma}{4},
 \sigma\sqrt{2\gamma\kappa_\delta}\right\},
 \qquad
 \varepsilon=\frac{c_*}{\sqrt\rho},
\]
and let \(Q\) keep the superior mean fixed while changing the inferior mean
to \(-\gamma/2-\varepsilon\).  We have already proved \(\rho\geq1\), so
\[
 \frac\gamma2+\varepsilon
 \leq\frac\gamma2+\frac{2M-\gamma}{4}<M.
\]
Hence all means under both laws lie in \([-M,M]\).  Their constant means are
H\"older for the declared \(L\), the uniform context density satisfies
\(c_X\leq1\leq C_X\), and the Gaussian model has the declared scale.  Their
constant gaps are respectively \(\gamma\) and \(\gamma+\varepsilon\), so the
fixed-gap and local-margin conditions hold by the preceding construction.
Their benchmark inferior-treatment mass equals \(b\geq b_0\).  Consequently,
\[
 P,Q\in\mathcal P_{\Delta_0,b_0}^{\mathrm{fg}}.
\]

Fix any
\(\mathsf D\in\mathfrak D_T(\rho;
\mathcal P_{\Delta_0,b_0}^{\mathrm{fg}})\), and define
\[
 N_{\mathrm{inf}}=\sum_{t=1}^T\mathbf1\{A_t=a_{\mathrm{inf}}\}.
\]
Under \(P\), each inferior-arm assignment has participant loss \(\gamma\).
Budget feasibility and the welfare-loss definition therefore give
\[
 \gamma\,\mathbb E_{P,\mathsf D}N_{\mathrm{inf}}
 =\mathcal R_T(P,\mathsf D)\leq\rho,
 \qquad
 \mathbb E_{P,\mathsf D}N_{\mathrm{inf}}\leq\frac\rho\gamma.
\]
The two laws have the same context distribution and differ only by
\(\varepsilon\) in the inferior-arm mean.  The exact adaptive Gaussian
likelihood identity in the coverage--information converse yields
\[
 \begin{aligned}
 \mathrm{KL}(\mathbb P_{P,\mathsf D}^{T}
             \Vert\mathbb P_{Q,\mathsf D}^{T})
 &=\frac{\varepsilon^2}{2\sigma^2}
   \mathbb E_{P,\mathsf D}N_{\mathrm{inf}}\\
 &\leq\frac{c_*^2}{2\sigma^2\gamma}
 \leq\kappa_\delta.
 \end{aligned}
\]
This includes the required inequality in the final comparison.

If arm one is superior, the targeting-gain definition gives
\(\theta(Q)-\theta(P)=\varepsilon
\int(1-\pi_0(x))\,dx\); if arm zero is superior, it gives
\(\theta(Q)-\theta(P)=\varepsilon\int\pi_0(x)\,dx\).  Hence in either case
\[
 |\theta(Q)-\theta(P)|=\varepsilon b\geq\frac\varepsilon2.
\]
For every
\(I_T\in\mathfrak C_{T,\delta}(\mathsf D;
\mathcal P_{\Delta_0,b_0}^{\mathrm{fg}})\), the established
coverage--information converse now implies
\[
 \sup_{R\in\mathcal P_{\Delta_0,b_0}^{\mathrm{fg}}}
 \mathbb E_{R,\mathsf D}|I_T(O_{1:T})|
 \geq
 \frac{1-2\delta}{2}|\theta(Q)-\theta(P)|
 \geq\frac{(1-2\delta)c_*}{4}\rho^{-1/2}.
\]
The design and honest interval were arbitrary.  Infimizing first over honest
intervals and then over budget-feasible non-anticipating designs proves the
lower bound with \(c=(1-2\delta)c_*/4\).  Together with the explicit upper
construction, this establishes both inequalities for all \(T\geq T_0\) and
\(C_1\log T\leq\rho\leq c_1T\).
